\section{\MakeUppercase {Research Directions for Digital Organizational Resilience}}
\label{sec:RD}
\renewcommand{\baselinestretch}{2.0}\normalsize


The Enactment--Scope framework provides an important stepping stone for the understanding of resilience, and its value grows further when it is used to generate concrete research questions. This section develops that next step. The focus here is digital organizational resilience, and especially cyber resilience, because digital environments make the core analytical problems particularly visible. In such settings, the substantial work already accomplished by regulators, standard-setting bodies, practitioners, and researchers can be extended by studying resilience as a socio-technical capability that is built, enacted, and tested under conditions of persistent disruption.\\

This shift extends the object of inquiry. When resilience is also examined through observable behavior, infrastructure, and governance arrangements, it becomes possible to evaluate it more directly and to refine the conceptual commitments practitioners and policymakers have already articulated. When it is treated as cumulative rather than static, questions of learning, memory, reinforcement, and erosion move to the center. And in digitally mediated environments, adaptive capacity is shaped not only by managerial judgment and organizational culture, but also by architectures, platforms, institutional settings, and the ways these elements condition what forms of response are possible.\\

The research agenda that follows is organized around two linked questions. The first concerns \emph{enactment}: through what mechanisms does resilience move from stated intention to operational capability? The second concerns \emph{scope}: at what organizational and institutional levels do these mechanisms operate? These questions build on the earlier diagnosis of resilience research. Where organizations have already developed indicators, governance structures, and evaluative benchmarks, the productive next step is to deepen the connection between observed adaptive performance and decision authority. At the same time, resilience in digitally mediated environments is distributed. It is not located in one team, one platform, or one formal rule. It is produced across routines, infrastructures, institutional arrangements, and inter-organizational ecosystems.\\

This is also where information systems research has a distinctive advantage. The field already addresses digital infrastructures, platform governance, socio-technical coordination, and cross-boundary information flows. These are not peripheral concerns here. They are the substance of the problem. The sections below therefore develop three research directions: operationalizing resilience enactment, examining resilience across organizational and institutional scopes, and tracing how resilience evolves over time.

\subsection{Operationalizing Resilience Enactment}

A first priority is to make resilience observable in practice. That is the central task of the enactment dimension of the framework. Organizations have invested significantly in resilience-related strategic language, governance documents, and policy discourse, and the next step is to develop equally robust ways of identifying resilience as it unfolds during disruption. IS research is well placed to support that step because it can connect behavioral indicators, computational traces, and organizational processes in ways that make adaptive capacity empirically visible.\\
\\
Building on the conceptual ground already established, research can now specify how resilience appears in practice and how it can be distinguished from related constructs such as preparedness or adaptability. Comparative work across domains already suggests that resilience emerges through the interaction of routines, infrastructures, and institutional processes \autocite{pumpuni2017resilience,manikis2019computational}. What remains to be developed is a more systematic basis for turning those observations into measures and design principles that organizations can readily apply.\\

Figure \ref{fig:triangle} supports that move by connecting theory, measurement, and governance as the three domains through which resilience becomes enacted capability. Each side of the triangle captures a different mechanism of translation. \emph{Operational indicators} render resilience empirically observable, complementing the conceptual commitments organizations already articulate. \emph{Design principles} allocate decision rights and structure adaptive responses. \emph{Benchmarks and warnings} align incentives while producing actionable signals. The figure matters less as a visual device than as an analytical claim: resilience becomes more readily governable when these three domains are linked rather than developed separately.\\

A converging line of work across several disciplines now develops behavioral and computational indicators that make resilience empirically assessable. Engineering resilience research has long offered formal candidates---recovery speed, reliability, vulnerability, and post-event functionality \autocite{hashimoto1982reliability,bruneau2003framework,francis2014metric,panteli2017metrics}. Public health has produced indicator suites for community and system resilience \autocite{cutter2010disaster,kruk2017building,haldane2021health}. More recent work in cybersecurity and digital risk management has begun to translate these traditions into organizational settings by characterizing investment--exposure dynamics, coordination performance, and incident-response patterns through quantitative and computational analysis \autocite{kuypers2018designing,gillard2023efficient,maillart2024computational,rane2024artificial}. Indicators such as coordination speed, the ability to reconfigure infrastructures, and the variability of response under pressure can complement the strategic commitments organizations have already made with directly measurable counterparts \autocite{mahajan2022participatory}. They also make comparison possible: organizations can benchmark themselves, and researchers can begin to accumulate evidence across cases. Importantly, developing such indicators is not solely a technical exercise. It also rests on a theoretical account of how behaviors at individual and collective levels scale into system-level outcomes \autocite{luthar2000construct,ungar2021modeling,kalisch2024neurobiology}.\\

Digital environments make this task especially urgent. Disruptions unfold quickly, interactions are distributed, and many adaptive signals are embedded in digital traces rather than retrospective narratives. Work in behavioral analytics and AI-assisted modeling therefore offers a useful path forward, particularly where resilience indicators can be calibrated against observable signatures in digital artifacts \autocite{rane2024artificial}.\\

Computational approaches, including agent-based modeling and simulation, are especially useful here. They allow researchers to explore how coordination patterns, incentive structures, and infrastructural configurations shape the capacity of organizations to absorb, reconfigure, and respond to disruption \autocite{arnold2004informationsharing,gillard2023efficient}. Used carefully, such methods do more than illustrate scenarios. They allow theory to be refined against simulated and empirical evidence and help bridge the gap between conceptual language and organization-level practice. In that sense, operationalization is not only about measurement. It is also about making resilience a capability that can be examined, compared, and deliberately improved.

\subsection{Resilience Across Organizational and Institutional Scopes}

A second research direction develops the scope dimension of the Enactment--Scope framework by asking where resilience mechanisms operate and how they are distributed across organizational, infrastructural, institutional, and inter-organizational settings.\\
\\
In digitally mediated environments, adaptive capacity is inherently distributed. It does not sit neatly inside a single team, technology, or governance layer. Governance architectures can enable resilience, but they can also constrain it by creating ambiguity in decision rights, accountability gaps, and escalation failures \autocite{ostrom1990governing,lebel2006governance}. These frictions may remain latent in stable periods, but they become visible under stress. They appear as delayed responses, constrained options, unclear ownership, and coordination breakdowns.\\

Studying these dynamics calls for approaches that connect institutional theory, regulatory ethnography, and actor-network perspectives with empirical work on digital infrastructures and everyday organizational practice. The point is to understand how resilience is enabled or constrained in the lived realities of digitally coordinated work, complementing the formal schemes of control already developed.\\

Table~\ref{fig:resilience_interface_opportunity} introduces an opportunity matrix that maps six research directions for digital organizational resilience---operational indicators, integrative frameworks, comparative and cumulative work, temporal and process markers, normative or value-based analysis, and governance and incentives---across four scopes: \emph{micro-level routines}, \emph{meso-level infrastructures}, \emph{macro-level institutions}, and \emph{inter-organizational ecosystems}. The cell shading indicates where each direction is currently most tractable and where it requires further institutional or methodological investment. The matrix does not replace the enactment perspective; it complements it. Figure~\ref{fig:triangle} clarifies \emph{how} resilience becomes operationalized through indicators, design principles, and evaluative benchmarks; Figure~\ref{fig:enactment_scope} clarifies \emph{where} those mechanisms operate; the opportunity matrix indicates, for each research direction, \emph{where the next round of evidence is most likely to land}.\\

\begin{table}
\small
\begin{tabular}{@{}>{\raggedright}p{4.2cm} *{4}{>{\centering\arraybackslash}p{2.2cm}}@{}}
\toprule
\textbf{Directions} & \textbf{Routines (micro)} & \textbf{Infrastructures (meso)} & \textbf{Institutions (macro)} & \textbf{Ecosystems (inter-org)} \\
\midrule
Operational indicators & \Closed & \Closed & \OpenClosed & \Closed \\
Integrative frameworks & \Open & \Closed & \Closed & \Closed \\
Comparative \& cumulative & \Open & \Closed & \OpenClosed & \Open \\
Temporal/process markers & \Closed & \Open & \Closed & \OpenClosed \\
Normative/value-based & \Open & \OpenClosed & \Closed & \Closed \\
Governance and incentives & \Open & \OpenClosed & \Closed & \Closed\\
\bottomrule
\end{tabular}
\vspace{2mm}
\footnotesize Legend: $\Closed$ high opportunity; $\OpenClosed$ promising ; $\Open$ low opportunity.
\caption{Opportunity matrix locating research directions across organizational scopes (micro routines, meso infrastructures, macro institutions, inter-organizational ecosystems).}
\label{fig:resilience_interface_opportunity}
\end{table}

Reading the matrix along its rows brings these opportunities into focus. \emph{Operational indicators} (row 1) are most tractable at the micro, meso, and ecosystem scopes, where behavioural traces, infrastructure telemetry, and information-sharing signals are increasingly available for empirical work \autocite{bruneau2003framework,francis2014metric,panteli2017metrics,kuypers2018designing,gillard2023efficient}. The macro scope is currently rated as promising rather than high because authoritative cross-jurisdictional indicators remain harder to construct, although the harmonized incident classification introduced by DORA and the Basel Committee's principles for operational resilience offer a useful starting point for comparable macro-level measurement \autocite{bcbs2021operational,european2022regulation,nist2024csf}. \emph{Comparative and cumulative} research (row 3) shows a similar profile: micro-level case comparison and ecosystem-level studies are already well supported by existing methods \autocite{barasa2018what,haldane2021health,maillart2024computational}, while macro-level comparability across jurisdictions is improving but still partial \autocite{meerow2016defining,allen2018quantifying}. \emph{Temporal and process markers} (row 4) are particularly tractable at the infrastructural and ecosystem scopes, where event logs, recovery curves, and longitudinal incident records support process tracing \autocite{ives1995measuring,leslie2013response,aase2020resilience,krakovska2023resilience}. \emph{Governance and incentives} (row 6) are most directly addressable at the meso scope, where infrastructure governance, access design, and escalation pathways are concrete enough to be studied empirically \autocite{ostrom1990governing,lebel2006governance,boin2013resilient,finma2023operational}. Read together, the three artefacts---triangle, matrix, and opportunity table---thus specify what to study, where the mechanisms operate, and where empirical investment is most likely to be productive.\\

A related issue concerns cooperation under pressure. Time-critical crises rarely respect organizational boundaries. Response depends on whether actors can coordinate quickly across firms, infrastructures, and institutional settings. Yet this is exactly where trust asymmetries, activation delays, and incentive misalignments often undermine performance. Theories of collective action, modular governance, and platform-mediated collaboration therefore provide an important starting point, especially in digital threat environments where rapid mobilization is essential \autocite{ostrom1990governing,brauer2008compartmental,mahajan2022participatory}.\\

Empirical work on modular security platforms such as MISP is especially relevant here because it shows how structural design can enable or inhibit timely collaboration. Participatory simulations and controlled experiments add another layer by identifying thresholds for cooperative engagement and the points at which collaboration begins to fail structurally \autocite{skopik2016problem,gillard2023efficient}. Disaster-response research offers additional design principles for cross-organizational information sharing through IT systems, and many of those insights are transferable to adversarial digital environments \autocite{arnold2004informationsharing}. The depth of resilience as a substantive capability depends in no small part on how these institutional arrangements and cooperation dynamics are configured across scopes.

\subsection{Temporal Dynamics of Resilience Enactment}

A third research direction extends the enactment perspective by asking how resilience is formed and re-formed over time. The point is not simply to add a temporal layer to the studies discussed above, but to examine how enacted capacities are built, stabilized, weakened, or redirected through experience, feedback, and institutional choice. On this view, resilience is neither a fixed organizational attribute nor something secured by resource availability alone. It is a contingent and relational capacity whose durability depends on whether earlier adaptations are reinforced, remembered, and carried forward, or instead neglected and allowed to erode \autocite{rutter1987psychosocial,luthar2000construct,ungar2021modeling}. This temporal dimension is comparatively underdeveloped in the IS literature, which still tends to focus on discrete events or isolated responses. Organizations do not confront each disruption from a blank slate: they carry residues of earlier crises---routines, assumptions, temporary fixes, and interpretive habits that, in \citeauthor{lanzara1999transient}'s analysis of post-disaster organizing, function as transient yet consequential scaffolds for later action \autocite{lanzara1999transient}. \textcite{westrum2004typology}'s typology of generative, bureaucratic, and pathological information cultures provides a complementary anchor for thinking about whether such residues become durable organizational learning or remain inert. Any persuasive account of resilience therefore has to explain not only how organizations respond in the moment, but also how adaptive capacity is retained, revised, or lost across successive episodes of disruption.\\
\\
Useful foundations for this line of inquiry already exist in systems theory and ecology. \citeauthor{holling1973resilience}'s work on adaptive cycles draws attention to recurring phases of growth, conservation, release, and reorganization, making clear that resilience is tied to sequences and transitions rather than to any single moment of adjustment \autocite{holling1973resilience,gunderson2002panarchy,folke2010resilience}. \citeauthor{wiener1948cybernetics}'s cybernetics places feedback at the center and shows how systems may be stabilized through corrective feedback or pushed toward instability through reinforcing feedback \autocite{wiener1948cybernetics}. More recent work on basin stability and return-time dynamics formalizes the same intuitions for nonlinear systems \autocite{gao2016universal,krakovska2023resilience}. Taken together, these traditions suggest that resilience depends not only on an organization's immediate ability to cope with disruption, but also on whether effective responses are reinforced, whether lessons become institutionalized, and whether failure produces revision rather than repetition \autocite{folke2004regime,allen2018quantifying}. Under some conditions resilience accumulates over time; under others it decays through forgetting, inertia, or maladaptation; and under the antifragile conditions described by \textcite{taleb2012antifragile}, controlled exposure to stressors may even compound adaptive capacity, provided the supporting learning routines remain durable.\\

Studying these dynamics calls for methodological pluralism. Longitudinal case studies are needed to follow how organizations adapt across repeated disruptions rather than single events. Process tracing can then identify the pathways through which learning practices, governance arrangements, and changes in technical systems shape later outcomes. Computational models offer complementary leverage by making it possible to examine how different rates of reinforcement, forgetting, or knowledge transfer generate different long-run resilience trajectories. Field evidence from Appendix~\ref{appendix:C} (theme iv) underscores why this temporal lens matters in practice: respondents repeatedly described post-incident reviews as conducted but rarely institutionalized (\textit{``The same issues resurface again and again because lessons learned are not retained''}), making the half-life of organizational learning a tractable and high-value object of empirical study. Taken together, these approaches would allow IS research to move beyond largely static conceptions of resilience and instead analyze it as a socio-technical capability that can strengthen, deteriorate, or mutate as organizations are repeatedly disturbed.

\subsection{Toward an Integrative Enactment-Scope Research Agenda}

Taken together, these research directions point toward a broader agenda: resilience can be studied most fruitfully by linking mechanisms of enactment to the organizational scopes within which they operate. Building on the contributions of fragmented perspectives, the field is now positioned to develop an integrated socio-technical account that relates behavioral processes, institutional conditions, and temporal development without collapsing them into a single level of analysis. This ambition is consistent with systems and ecological traditions that treat resilience as an emergent effect of feedback, interdependence, and adaptation \autocite{wiener1948cybernetics,holling1973resilience}. For IS research, that means building explanations that connect micro-level action, meso-level organizational and infrastructural dynamics, and macro-level governance constraints. It also means developing indicators that complement compliance reporting with finer-grained traces of adaptation as it unfolds: changes in response variability, shifts in coordination under stress, and the reconfiguration of technical and organizational arrangements during disruption \autocite{aerts2018integrating}. Anchoring resilience in observable behavior is the bridge that helps the conceptual ambition of the field meet analytical precision.\\
\\
At the same time, resilience does not emerge in a vacuum. The institutional and organizational environments in which actors operate shape what forms of adaptation are feasible in the first place. Rules, incentives, authority structures, and escalation pathways do not simply constrain action after disruption; they structure the range of possible responses before disruption occurs \autocite{ostrom1990governing,ungar2021modeling}. Existing field-based and computational studies already show that these conditions matter, often appearing as friction, delay, coordination breakdown, or reduced strategic flexibility. The harder task now is to turn these scattered insights into analytical frameworks that can travel across settings without losing explanatory force.\\

Figure \ref{fig:resilience_interface_opportunity} makes this research space visible by locating opportunities across micro behaviors, meso infrastructures, macro institutions, and inter-organizational ecosystems. Read together, Figures \ref{fig:triangle} and \ref{fig:resilience_interface_opportunity} define the central logic of the Enactment--Scope framework. The triangle specifies how resilience moves from articulated organizational commitment to enacted practice. The opportunity matrix shows where those mechanisms can be studied and where intervention is most likely to matter. Advancing this agenda will require research programs that combine theoretical ambition with methodological range, including computational modeling, longitudinal case analysis, and field-based validation. Used together, these approaches offer a more credible path toward cumulative and transferable knowledge about organizational resilience \autocite{teece2007explicating,gillard2023efficient,kalisch2024neurobiology}. More importantly, they provide a stronger basis for IS theorizing about digital resilience by showing how adaptive capacity emerges through the interaction of technological architectures, organizational practices, and institutional environments. On this view, resilience is best understood as more than a strategic statement or a checklist: it is an enacted, distributed, and revisable socio-technical capability for maintaining and reorganizing coordinated action under persistent uncertainty \autocite{weick2011managing,woods2015four}.